\documentclass[11pt]{article}
% Shared preamble for the rigor documents (Lamport-style structured proofs).  \input this file after \documentclass.
\usepackage[T1]{fontenc}\usepackage{lmodern}\usepackage{microtype}
\usepackage[margin=1in]{geometry}
\usepackage{amsmath,amssymb,amsthm,mathtools}
\usepackage{xcolor}\usepackage{enumitem}\usepackage{booktabs}\usepackage{graphicx}\usepackage{hyperref}
\usepackage[most]{tcolorbox}
\definecolor{navy}{HTML}{17365D}\definecolor{teal}{HTML}{117A8B}\definecolor{warn}{HTML}{A33A2B}\definecolor{okgreen}{HTML}{2E7D4F}
\hypersetup{colorlinks=true,linkcolor=navy,citecolor=teal,urlcolor=teal}
\theoremstyle{plain}
\newtheorem{theorem}{Theorem}[section]\newtheorem{lemma}[theorem]{Lemma}\newtheorem{proposition}[theorem]{Proposition}\newtheorem{corollary}[theorem]{Corollary}
\theoremstyle{definition}\newtheorem{definition}[theorem]{Definition}\newtheorem{remark}[theorem]{Remark}\newtheorem{claim}[theorem]{Claim}\newtheorem{issue}[theorem]{Issue}
% ---------- Lamport structured proofs ----------
% \lstep{level}{number}{assertion}   -> indented "<level>number. assertion"
% \lproof                             -> "PROOF:" line (start of the sub-proof of the preceding step)
% \lby{justification}                 -> "BY ..." leaf justification for the preceding step
% \lqed{level}{number}                -> QED step of a sub-proof
% keywords: \lassume \lprove \lsuffices \lcase \ldefine \lpick \lnew
\newlength{\lindent}\setlength{\lindent}{1.7em}
\newcommand{\lstep}[3]{\par\smallskip\noindent\hspace*{\dimexpr\numexpr#1-1\relax\lindent\relax}{\color{navy}$\langle#1\rangle#2$.}\ #3}
\newcommand{\lproof}{\par\noindent\hspace*{\lindent}{\scshape Proof:}}
\newcommand{\lby}[1]{\par\noindent\hspace*{\lindent}{\scshape By}\ #1.}
\newcommand{\lqed}[2]{\par\smallskip\noindent\hspace*{\dimexpr\numexpr#1-1\relax\lindent\relax}{\color{navy}$\langle#1\rangle#2$.}\ {\scshape Q.E.D.}}
\newcommand{\lassume}{{\scshape Assume}\ }\newcommand{\lprove}{{\scshape Prove}\ }\newcommand{\lsuffices}{{\scshape Suffices}\ }
\newcommand{\lcase}{{\scshape Case}\ }\newcommand{\ldefine}{{\scshape Define}\ }\newcommand{\lpick}{{\scshape Pick}\ }\newcommand{\lnew}{{\scshape New}\ }
% ---------- verdict boxes ----------
\newtcolorbox{verdictok}{colback=okgreen!8,colframe=okgreen,title={\bfseries Verdict: verified},fonttitle=\small}
\newtcolorbox{verdictgap}{colback=orange!10,colframe=orange!80!black,title={\bfseries Verdict: gap / caveat},fonttitle=\small}
\newtcolorbox{verdictbreak}{colback=warn!10,colframe=warn,title={\bfseries Verdict: proof-breaking issue},fonttitle=\small}
\newtcolorbox{citedbox}[1]{colback=navy!5,colframe=navy!60,title={\bfseries Cited result: #1},fonttitle=\small,breakable}
\setlength{\parindent}{0pt}\setlength{\parskip}{4pt}\allowdisplaybreaks
\newcommand{\cA}{\mathcal A}\newcommand{\cF}{\mathcal F}\newcommand{\cM}{\mathfrak M}\newcommand{\cN}{\mathfrak N}

% extra calligraphic shortcuts (\providecommand: no clash if a document defines its own)
\providecommand{\cB}{\mathcal B}\providecommand{\cC}{\mathcal C}\providecommand{\cD}{\mathcal D}\providecommand{\cE}{\mathcal E}\providecommand{\cG}{\mathcal G}\providecommand{\cH}{\mathcal H}\providecommand{\cI}{\mathcal I}\providecommand{\cJ}{\mathcal J}\providecommand{\cK}{\mathcal K}\providecommand{\cL}{\mathcal L}\providecommand{\cO}{\mathcal O}\providecommand{\cP}{\mathcal P}\providecommand{\cQ}{\mathcal Q}\providecommand{\cR}{\mathcal R}\providecommand{\cS}{\mathcal S}\providecommand{\cT}{\mathcal T}\providecommand{\cU}{\mathcal U}\providecommand{\cV}{\mathcal V}\providecommand{\cW}{\mathcal W}\providecommand{\cX}{\mathcal X}\providecommand{\cY}{\mathcal Y}\providecommand{\cZ}{\mathcal Z}
\newcommand{\Fid}{\operatorname{F}}\newcommand{\Tr}{\operatorname{Tr}}\newcommand{\eps}{\varepsilon}\newcommand{\dd}{\,\mathrm d}

\providecommand{\one}{\mathbf 1}\providecommand{\Ran}{\operatorname{Ran}}\providecommand{\norm}[1]{\lVert #1\rVert}
\newcommand{\Nst}[1]{\lVert #1\rVert_{\star}}\newcommand{\Nsto}[1]{\lVert #1\rVert_{\star\star}}\newcommand{\Pinch}{\mathcal P}

\newcommand{\Qt}{\widetilde{Q}}
\newcommand{\Xo}{X_{11}}
\newcommand{\Xt}{X_{22}}
\newcommand{\Xc}{X_{12}}
\newcommand{\HS}{\mathfrak{S}_2}
\newcommand{\kc}{\kappa_c}
\newcommand{\zt}{\zeta}
\newcommand{\Phiw}{\Phi_W}

\title{Certified numerics for $\Phi(\zeta)$:\\
a quadratic block bound for $\sqrt{\Qt}$ and ball-arithmetic enclosures}
\author{Agent RC (Phase 2b, track RC)}
\date{2026-09-09}

\begin{document}
\maketitle
\tableofcontents

\section{Scope and summary}
\label{sec:scope}
% filled in last
This note has two independent parts, both aimed at the width of the error brackets for
$\Phi(\zeta)=-\log\Fid(\omega,\widetilde\omega_s)$ of Note~5, Table~4 (\texttt{tab:scan}; previously cited here as Table~3), which
\texttt{tail\_bound.tex} (agent PT) left at $82$--$154\,\%$, with bars that are not certified: their upper half rests on the numerical inputs of the final bound of \texttt{tail\_bound.tex} (see Correction~C5).

\paragraph{Part I (analysis, \S\ref{sec:sqrt}--\S\ref{sec:nocross}).}  The open item of
\texttt{tail\_bound.tex}, ``a bound on $\|\sqrt{\widetilde Q}-\sqrt{X_{11}\oplus X_{22}}\|_2$
quadratic in the off-diagonal block $X_{12}$'', is settled as follows.
\begin{itemize}[leftmargin=1.6em]
\item The bound is \emph{false} in the unweighted form $\|\sqrt{\widetilde Q}-\sqrt{X_{11}\oplus X_{22}}\|_2^2
 \le C\|X_{12}\|_2^2$ with an absolute constant $C$: Proposition~\ref{prop:cex} exhibits
 $2\times2$ positive contractions with ratio $\ge3/(2\eps)$, $\eps=\lambda_{\min}(X_{11})$.
\item It is \emph{true} with the natural Daleckii--Krein weight.  Theorem~\ref{thm:sqrt}:
 with $a_i$ the eigenvalues of $X_{11}$, $b_j$ those of $X_{22}$ and $V_{ij}$ the matrix elements of
 $X_{12}$ in the corresponding eigenbases,
 \[
  \Bigl\|\sqrt{\widetilde Q}-\sqrt{X_{11}\oplus X_{22}}\Bigr\|_2^2
  =\sum_{\alpha,i}\frac{|\langle\phi_\alpha,X_{12}\psi_i\rangle+\langle\phi_\alpha,X_{12}^*\psi_i\rangle|^2}
  {(\sqrt{\beta_\alpha}+\sqrt{a_i})^{2}},
 \]
 an \emph{identity} (Theorem~\ref{thm:sqrt}) with $\beta_\alpha$, $a_i$ the eigenvalues of
 $\widetilde Q$ and of its pinching: the Daleckii--Krein weight $(\sqrt{\beta}+\sqrt a)^{-2}$ is
 exactly what replaces the missing Lipschitz constant.
\item The weighted norm is \emph{finite and small} for the actual kernel: positivity alone gives
 (Corollary~\ref{cor:schur}) $\Nst{X_{12}}^2\le\tau_++\tau_-=\tau$ --- the same occupation tail that
 governs everything else, so the weight costs nothing.  This is the exact point at which the
 Schur-complement inequality $X_{12}^*X_{11}^{-1}X_{12}\le X_{22}$, guessed by PT to be ``the right
 tool'', does the work.
\item Consequently (Theorem~\ref{thm:nocross}) the ultraviolet tail bound can be freed of the Bures
 cross term: with $\Phi_{22}$ the complement-block defect and $\Phi_{\rm off}=-\log|\det(1+T)|$,
 \[
  0\ \le\ \Phi-\Phi_W\ \le\ \Phi_{22}+\Phi_{\rm off},
 \]
 an inequality between \emph{computable} quantities with no term of order $\sqrt\Phi$.  The
 mechanism is that the optimal Uhlmann purification of the pair $(\omega_Q,\omega_A)$,
 $A=\Pinch\widetilde Q$, can be chosen \emph{block diagonal}, so that the off-diagonal block enters
 $\log\det$ only at second order.  No closed-form constant $C_\star$ is claimed: the a priori route
 through hypothesis \eqref{eq:Hyp} has a measured constant $c_0\in[18,42]$ and is about four to nine times better than the $82$--$154\,\%$ of \texttt{tail\_bound.tex} but
 much weaker than the computable form: $17.8$--$19.9\,\%$ against $0.708$--$1.010\,\%$ of $\Phi_W$ at $\kappa_c=25.3$ (\S\ref{sec:num}, \S\ref{sec:rcr}; see Correction~C2).
\end{itemize}
Numerically this replaces the factor-$80$ overshoot of \texttt{tail\_bound.tex} by a factor
$116$--$153$ at $\kappa_c=25.3$ (Table~\ref{tab:rcr}; the $50$--$160$ printed here before used the withdrawn double-precision widths); the resulting widths are stated in \S\ref{sec:rcr}, which supersedes
Table~\ref{tab:cert} after the referee report.

\paragraph{Part II (arithmetic, \S\ref{sec:ball}--\S\ref{sec:table}).}  The $3\,\%$ arithmetic margin
of \texttt{tail\_bound.tex}, \S8, is replaced by a rigorous \texttt{arb}/\texttt{acb} ball-arithmetic
evaluation of the whole windowed pipeline: \emph{delivered}: (a) the box symbol $Q_N$ enclosed
entrywise to $5.62\times10^{-18}$ from $2N$ certified one-dimensional integrals, with no
eigendecomposition of a ball matrix anywhere (\S\ref{sec:qbox}); (b) the downstream fidelity
quantities in \texttt{mpmath} at $60$ digits from that enclosure (\S\ref{sec:rcr}).
\emph{Not delivered, and stated as such}: a ball-arithmetic evaluation of the defect $\widehat D$
(\S\ref{sec:dhat}), verified Rayleigh--Ritz eigenvalue enclosures, and a ball-arithmetic
determinant; \S\ref{sec:eig} records the recipes only.  \S\ref{sec:rcr} states exactly what is
certified and supersedes Table~\ref{tab:cert}.
\section{The square root of a block perturbation}\label{sec:sqrt}

Throughout this section $\mathfrak h=\mathfrak h_1\oplus\mathfrak h_2$ is finite dimensional,
\[
 B=\begin{pmatrix}A_1&W\\ W^*&A_2\end{pmatrix}\ \ge\ 0,\qquad
 A:=A_1\oplus A_2=\Pinch(B),\qquad V:=B-A=\begin{pmatrix}0&W\\ W^*&0\end{pmatrix},
\]
where $\Pinch$ is the pinching (conditional expectation onto the block-diagonal subalgebra).  We fix
orthonormal eigenbases $A_1u_i=a_iu_i$, $A_2w_j=b_jw_j$ and write $V_{ij}=\langle u_i,Ww_j\rangle$.
Define the weighted quantities
\begin{equation}\label{eq:starnorm}
 \Nst{W}^2:=\sum_{i,j}\frac{|V_{ij}|^2}{\bigl(\sqrt{a_i}+\sqrt{b_j}\bigr)^2},
 \qquad
 \Nsto{W}:=\sum_{i,j}\frac{|V_{ij}|^2}{\sqrt{a_ib_j}\,\bigl(\sqrt{a_i}+\sqrt{b_j}\bigr)} .
\end{equation}
(The second is used only in the remainder estimate; both are $+\infty$ if some $a_i$ or $b_j$
vanishes with $V_{ij}\neq0$, which positivity of $B$ forbids, see Lemma~\ref{lem:2x2}.)

\begin{lemma}[$2\times2$ positivity]\label{lem:2x2}
$|V_{ij}|^2\le a_ib_j$ for all $i,j$.
\end{lemma}
\lstep{1}{1}{The $2\times2$ principal submatrix of $B$ in the orthonormal system $(u_i,w_j)$ is
$\begin{pmatrix}a_i&V_{ij}\\ \overline{V_{ij}}&b_j\end{pmatrix}$.}
\lby{$\langle u_i,Bu_i\rangle=\langle u_i,A_1u_i\rangle=a_i$ (the $(1,2)$ block of $B$ does not
contribute to a diagonal $\mathfrak h_1$ entry), likewise $\langle w_j,Bw_j\rangle=b_j$, and
$\langle u_i,Bw_j\rangle=\langle u_i,Ww_j\rangle=V_{ij}$}
\lstep{1}{2}{A principal submatrix of a positive semidefinite matrix is positive semidefinite, and a
positive semidefinite $2\times2$ matrix has nonnegative determinant.}
\lby{$\langle f,Bf\rangle\ge0$ for $f$ in the span of $u_i,w_j$; and $\det\ge0$ is the product of the
two nonnegative eigenvalues}
\lqed{1}{3}\lby{$\langle1\rangle1$, $\langle1\rangle2$: $a_ib_j-|V_{ij}|^2\ge0$}

\subsection{The exact Sylvester representation and the sharp weight}\label{ss:sylv}

\begin{theorem}[exact form of the square-root difference]\label{thm:sqrt}
Let $A,B\ge0$ on a finite-dimensional space, $V:=B-A$, and let $A\psi_i=a_i\psi_i$,
$B\phi_\alpha=\beta_\alpha\phi_\alpha$ be orthonormal eigendecompositions.  Then
\begin{equation}\label{eq:sylv}
 B-A=\sqrt B\,\Xi+\Xi\,\sqrt A,\qquad \Xi:=\sqrt B-\sqrt A,
\end{equation}
and, if $A$ and $B$ are invertible,
\begin{equation}\label{eq:xiexact}
 \boxed{\ \bigl\|\sqrt B-\sqrt A\bigr\|_2^2
 \;=\;\sum_{\alpha,i}\frac{\bigl|\langle\phi_\alpha,V\psi_i\rangle\bigr|^2}
 {\bigl(\sqrt{\beta_\alpha}+\sqrt{a_i}\bigr)^{2}}\ }
\end{equation}
Moreover, if $B=A+V$ with $A=\Pinch(B)$ block diagonal and $V$ purely off-diagonal, then the
first-order term
\begin{equation}\label{eq:firstorder}
 \mathcal L_A(V):=\frac1\pi\int_0^\infty t^{1/2}\frac1{A+t}\,V\,\frac1{A+t}\,\dd t,
 \qquad \bigl(\mathcal L_A(V)\bigr)_{ij}=\frac{V_{ij}}{\sqrt{a_i}+\sqrt{a_j}},
\end{equation}
is again purely off-diagonal, $\Pinch\mathcal L_A(V)=0$, and $\|\mathcal L_A(V)\|_2=\sqrt2\,\Nst{W}$
with $\Nst\cdot$ as in \eqref{eq:starnorm}.
\end{theorem}
\lstep{1}{1}{$\sqrt B\,\Xi+\Xi\sqrt A=\sqrt B\sqrt B-\sqrt B\sqrt A+\sqrt B\sqrt A-\sqrt A\sqrt A=B-A$.}
\lby{expansion of $\Xi=\sqrt B-\sqrt A$ and $(\sqrt B)^2=B$, $(\sqrt A)^2=A$}
\lstep{1}{2}{In the (mixed) bases $(\phi_\alpha)$, $(\psi_i)$, \eqref{eq:sylv} reads
$\langle\phi_\alpha,V\psi_i\rangle=(\sqrt{\beta_\alpha}+\sqrt{a_i})\langle\phi_\alpha,\Xi\psi_i\rangle$.}
\lby{$\sqrt B\phi_\alpha=\sqrt{\beta_\alpha}\phi_\alpha$, $\sqrt A\psi_i=\sqrt{a_i}\psi_i$, and
$\langle1\rangle1$}
\lstep{1}{3}{$\|\Xi\|_2^2=\sum_{\alpha i}|\langle\phi_\alpha,\Xi\psi_i\rangle|^2$, which with
$\langle1\rangle2$ gives \eqref{eq:xiexact}; the denominators are $>0$ because $A,B>0$.}
\lby{$\|\Xi\|_2^2=\Tr[\Xi^*\Xi]$ evaluated by inserting the two complete orthonormal systems}
\lstep{1}{4}{$\sqrt X=\frac1\pi\int_0^\infty X(X+t)^{-1}t^{-1/2}\dd t$ for $X>0$, hence
$\frac{\dd}{\dd\varepsilon}\big|_0\sqrt{A+\varepsilon V}=\mathcal L_A(V)$ and
$(\mathcal L_A(V))_{ij}=V_{ij}\cdot\frac1\pi\int_0^\infty\frac{t^{1/2}\dd t}{(a_i+t)(a_j+t)}
=\frac{V_{ij}}{\sqrt{a_i}+\sqrt{a_j}}$.}
\lby{the scalar identity $\frac1\pi\int_0^\infty\frac{x}{x+t}t^{-1/2}\dd t=\sqrt x$ (substitute
$t=xu^2$ and $\int_0^\infty\frac{2\dd u}{1+u^2}=\pi$), applied to the spectral measure of $X$;
differentiation under the integral sign is legitimate because the integrand is a rational function of
$\varepsilon$ with poles off $[0,\infty)$ for $|\varepsilon|$ small and is dominated by
$C\min(t^{-1/2},t^{-3/2})$; the resolvent identity
$\partial_\varepsilon(A+\varepsilon V+t)^{-1}=-(A+t)^{-1}V(A+t)^{-1}$ at $\varepsilon=0$; and
$\frac1\pi\int_0^\infty\frac{t^{1/2}}{(a+t)(b+t)}\dd t=\frac1{\sqrt a+\sqrt b}$, which follows from
the same substitution after partial fractions}
\lstep{1}{5}{If $V$ is off-diagonal for the decomposition $\mathfrak h_1\oplus\mathfrak h_2$ and $A$
is block diagonal, then $\mathcal L_A(V)$ is off-diagonal, so $\Pinch\mathcal L_A(V)=0$; and
$\|\mathcal L_A(V)\|_2^2=2\sum_{i,j}\frac{|V_{ij}|^2}{(\sqrt{a_i}+\sqrt{b_j})^2}=2\Nst{W}^2$.}
\lby{$(A+t)^{-1}$ is block diagonal, so $(A+t)^{-1}V(A+t)^{-1}$ is off-diagonal;
by $\langle1\rangle4$ the $(1,2)$ block has entries $V_{ij}/(\sqrt{a_i}+\sqrt{b_j})$ and the $(2,1)$
block is its adjoint, whence the factor $2$}
\lqed{1}{6}\lby{$\langle1\rangle1$--$\langle1\rangle5$}

\begin{remark}[why the weight is the right one]
\eqref{eq:xiexact} is an identity, not an estimate: no bound of the form
$\|\Xi\|_2\le c\,\|V\|_2$ can be better than the one obtained by replacing every
$(\sqrt{\beta_\alpha}+\sqrt{a_i})^{-1}$ by its maximum $\lambda_{\min}^{-1/2}$.  Discarding
$\sqrt{\beta_\alpha}$ gives the one-sided bound
\begin{equation}\label{eq:onesided}
 \|\Xi\|_2^2\ \le\ \Tr\bigl[V^*A^{-1}V\bigr]
 \ =\ \Tr\bigl[W^*A_1^{-1}W\bigr]+\Tr\bigl[WA_2^{-1}W^*\bigr]
\end{equation}
(and the same with $A$ replaced by $B$), which is what we use below.
\end{remark}

\begin{proposition}[the unweighted bound is false]\label{prop:cex}
For every $0<\delta<\eps<1/2$ the matrices
\[
 B=\begin{pmatrix}\eps&\sqrt{\eps\delta}\\ \sqrt{\eps\delta}&\delta\end{pmatrix}\ \ge0,
 \qquad A=\Pinch(B)=\begin{pmatrix}\eps&0\\0&\delta\end{pmatrix},\qquad W=\sqrt{\eps\delta},
\]
satisfy $0\le B\le1$, $\|W\|_2^2=\eps\delta$ and
\[
 \bigl\|\sqrt B-\sqrt A\bigr\|_2^2
 =\Bigl(\tfrac{\eps}{\sqrt{\eps+\delta}}-\sqrt\eps\Bigr)^2
 +\Bigl(\tfrac{\delta}{\sqrt{\eps+\delta}}-\sqrt\delta\Bigr)^2+\frac{2\eps\delta}{\eps+\delta}
 \ \xrightarrow[\ \delta/\eps\to0\ ]{}\ 3\delta\,(1+o(1)),
\]
so that $\|\sqrt B-\sqrt A\|_2^2\big/\|W\|_2^2\to3/\eps$.  Hence no inequality
$\|\sqrt B-\sqrt A\|_2^2\le C\|X_{12}\|_2^2$ with an absolute constant $C$ can hold for positive
contractions, and the blow-up rate as the window edge $\eps=\lambda_{\min}(X_{11})$ tends to $0$ is
exactly $\eps^{-1}$.
\end{proposition}
\lstep{1}{1}{$B=vv^*$ with $v=(\sqrt\eps,\sqrt\delta)^{\!\top}$, so $B\ge0$, $\|v\|^2=\eps+\delta<1$
and $B\le\|v\|^2\le1$.}
\lby{rank-one positivity; $\|B\|=\|v\|^2$}
\lstep{1}{2}{$\sqrt B=\|v\|^{-1}vv^*=(\eps+\delta)^{-1/2}B$.}
\lby{$B^2=\|v\|^2B$, so $(\|v\|^{-1}B)^2=B$ and $\|v\|^{-1}B\ge0$}
\lstep{1}{3}{The displayed formula for $\|\sqrt B-\sqrt A\|_2^2$, and its limit.}
\lby{$\sqrt A=\mathrm{diag}(\sqrt\eps,\sqrt\delta)$, $\langle1\rangle2$, and the entrywise
Hilbert--Schmidt norm; as $\delta/\eps\to0$ the three terms are
$\frac{\delta^2}{4\eps}(1+o(1))$, $\delta(1+o(1))$, $2\delta(1+o(1))$}
\lqed{1}{4}\lby{$\langle1\rangle1$--$\langle1\rangle3$ and $\|W\|_2^2=\eps\delta$}

\begin{remark}
In Proposition~\ref{prop:cex} the exact weighted expression \eqref{eq:xiexact} evaluates to
$\Nst{W}^2=\eps\delta/(\sqrt\eps+\sqrt\delta)^2=\delta(1+o(1))$: the weighted norm sees the correct
size $\delta$, and the extra factor $3$ comes from the eigenvalue $\beta=0$ of $B$, i.e.\ from the
\emph{non-perturbative} regime $\|W\|=\sqrt{\eps\delta}$ comparable to $\lambda_{\min}$.  This is the
precise sense in which the answer to PT's question is ``yes, with the Daleckii--Krein weight, and
never without it''.
\end{remark}

\subsection{Positivity makes the weighted norm small}\label{ss:schur}

\begin{corollary}[Schur-complement bound]\label{cor:schur}
Let $0\le\widetilde Q\le1$ (both inequalities are used), $B=\widetilde Q$, $A_1=X_{11}$,
$A_2=X_{22}$, $W^\pm:=E\widetilde QE^\pm$, and assume
\begin{equation}\label{eq:commhyp}
 [X_{22},E^+]=[X_{22},E^-]=0,
\end{equation}
so that the eigenbasis $(w_j)$ of $X_{22}$ splits into an $E^+$ and an $E^-$ family.  Then
\begin{equation}\label{eq:schurb}
 \Tr\bigl[(W^+)^*X_{11}^{-1}W^+\bigr]\le\tau_+,\qquad
 \Tr\bigl[(W^-)^*(1-X_{11})^{-1}W^-\bigr]\le\tau_-,
\end{equation}
and, provided $\tau_-<1$,
\begin{equation}\label{eq:starsmall}
 \Nst{X_{12}}^2\ \le\ \tau_++\frac{\|W^-\|_2^2}{1-\tau_-}\ \le\ \tau_++\frac{\tau_-}{1-\tau_-} .
\end{equation}
Without \eqref{eq:commhyp} the same proof gives \eqref{eq:starsmall} with $\tau_\pm$ replaced by the
traces of $X_{22}$, resp.\ $1-X_{22}$, over the spectral subspaces of $X_{22}$ on which its eigenvalue
is $\le\frac12$, resp.\ $>\frac12$.
\end{corollary}
\lstep{1}{1}{$X_{12}^*X_{11}^{-1}X_{12}\le X_{22}$ and $X_{12}^*(1-X_{11})^{-1}X_{12}\le1-X_{22}$.}
\lby{$\widetilde Q\ge0$ and $X_{11}>0$ give the Schur complement
$X_{22}-X_{12}^*X_{11}^{-1}X_{12}\ge0$ (factorise
$\widetilde Q=\binom{1\ \ 0}{X_{12}^*X_{11}^{-1}\ \ 1}\bigl(X_{11}\oplus(X_{22}-X_{12}^*X_{11}^{-1}X_{12})\bigr)
\binom{1\ \ X_{11}^{-1}X_{12}}{0\ \ 1}$); applying the same to $1-\widetilde Q\ge0$, whose blocks are
$1-X_{11}$, $-X_{12}$, $1-X_{22}$, gives the second.  \emph{Both} $\widetilde Q\ge0$ and
$\widetilde Q\le1$ are needed}
\lstep{1}{2}{Compressing the two inequalities of $\langle1\rangle1$ with $E^+$, resp.\ $E^-$, and
taking traces gives \eqref{eq:schurb}.}
\lby{$\Tr[E^+X_{22}E^+]=\tau_+$ and $\Tr[E^-(1-X_{22})E^-]=\tau_-$ by \eqref{eq:taudef} of
\texttt{tail\_bound.tex}; $Z\le Z'\Rightarrow\Tr[P^*ZP]\le\Tr[P^*Z'P]$}
\lstep{1}{3}{Under \eqref{eq:commhyp} every eigenvalue $b_j$ of $X_{22}$ belonging to the $E^-$
family satisfies $1-b_j\le\tau_-$, hence $b_j\ge1-\tau_-$.}
\lby{the $E^-$ family spans $E^-H$, on which $1-X_{22}$ is positive with trace $\tau_-$, so each of
its eigenvalues is at most $\tau_-$; \eqref{eq:commhyp} is what makes the $1-b_j$ eigenvalues of
$E^-(1-X_{22})E^-$}
\lstep{1}{4}{$\Nst{X_{12}}^2=\sum_{i,\,j\in E^+}\frac{|V_{ij}|^2}{(\sqrt{a_i}+\sqrt{b_j})^2}
+\sum_{i,\,j\in E^-}\frac{|V_{ij}|^2}{(\sqrt{a_i}+\sqrt{b_j})^2}
\le\Tr[(W^+)^*X_{11}^{-1}W^+]+\frac{\|W^-\|_2^2}{1-\tau_-}$.}
\lby{\eqref{eq:commhyp} lets the $j$-sum be split; on the $E^+$ part drop $\sqrt{b_j}$, giving the
weight $a_i^{-1}$ and the first trace; on the $E^-$ part $(\sqrt{a_i}+\sqrt{b_j})^2\ge b_j\ge1-\tau_-$
by $\langle1\rangle3$, so the weight is at most $(1-\tau_-)^{-1}$ --- \emph{not} $1$}
\lstep{1}{5}{$\|W^-\|_2^2\le\tau_-$.}
\lby{$\|EXE^-\|_2^2\le\|1-X_{11}\|\cdot\Tr[E^-(1-X)E^-]\le\tau_-$, the particle--hole mirror of
\texttt{tail\_bound.tex}, Theorem~4.1$\langle2\rangle1$, which again uses $\widetilde Q\le1$}
\lqed{1}{6}\lby{$\langle1\rangle2$, $\langle1\rangle4$, $\langle1\rangle5$}

\begin{remark}[status]
In the free-fermion application $E^\pm$ are spectral projections of $K=\log\frac{1-Q}Q$ and
$X_{22}=E^\perp\widetilde QE^\perp$ does \emph{not} commute with them exactly, so
\eqref{eq:commhyp} is an idealisation and the last sentence of the corollary is the applicable form.
Numerically (\S\ref{sec:num}) the ratio $\Nst{X_{12}}^2/\tau$ is $0.71$, so the conclusion holds
comfortably.
\end{remark}

\section{A tail bound without the Bures cross term}\label{sec:nocross}

The factor-$80$ overshoot of \texttt{tail\_bound.tex} comes from the term $2\sqrt{\Phi\Phi_\partial}$
produced by the Bures-angle triangle inequality.  We now remove it.  The mechanism is that the
\emph{optimal} Uhlmann purification of the pair $(\omega_Q,\omega_A)$, $A=\Pinch(\widetilde Q)$, is
block diagonal, and that a block-diagonal operator does not see the off-diagonal first-order
displacement, Theorem~\ref{thm:sqrt}$\langle1\rangle5$.

\begin{lemma}[pinching lowers the square root]\label{lem:jensen}
Let $\Pinch$ be the pinching onto the block diagonal of $\mathfrak h_1\oplus\mathfrak h_2$,
$B\ge0$, $A:=\Pinch(B)$, $\Xi_0:=\sqrt B-\sqrt A$ and $\Theta_0:=-\Pinch(\Xi_0)=\sqrt A-\Pinch\sqrt B$.
Then
\begin{equation}\label{eq:jensen}
 \Theta_0\ \ge\ 0,\qquad \Tr\bigl[\sqrt A\,\Theta_0\bigr]=\tfrac12\bigl\|\Xi_0\bigr\|_2^2 ,
\end{equation}
and for every Hermitian block-diagonal $Y$ with $-c\sqrt A\le Y\le c\sqrt A$,
\begin{equation}\label{eq:Ytrace}
 \bigl|\Tr[Y\,\Xi_0]\bigr|\ \le\ \tfrac{c}{2}\,\|\Xi_0\|_2^2 .
\end{equation}
\end{lemma}
\lstep{1}{1}{$\Pinch$ is a unital, trace-preserving, completely positive map with
$\Pinch\circ\Pinch=\Pinch$, and $\Pinch(Z)=Z$ for block-diagonal $Z$, $\Tr[Z\Pinch(Y)]=\Tr[\Pinch(Z)Y]$.}
\lby{$\Pinch(X)=\frac12(X+\Gamma X\Gamma)$ with $\Gamma=1_{\mathfrak h_1}\oplus(-1_{\mathfrak h_2})$
self-adjoint unitary; averages of unitary conjugations are unital completely positive and trace
preserving, and self-adjoint for the Hilbert--Schmidt pairing}
\lstep{1}{2}{$\Pinch(\sqrt B)\le\sqrt{\Pinch(B)}=\sqrt A$, i.e.\ $\Theta_0\ge0$.}
\lby{Jensen's operator inequality for the operator concave function $t\mapsto\sqrt t$ on $[0,\infty)$
and the unital positive map $\Pinch$: $\Pinch(f(B))\le f(\Pinch(B))$ (Hansen--Pedersen 2003,
Theorem~2.1; for a pinching this is Davis' theorem, and it also follows from
$\Pinch(X)=\frac12(X+\Gamma X\Gamma)$ together with the operator concavity of $\sqrt\cdot$,
$\frac12(\sqrt{B}+\sqrt{\Gamma B\Gamma})\le\sqrt{\frac12(B+\Gamma B\Gamma)}$)}
\lstep{1}{3}{$\Tr[\sqrt A\,\Theta_0]=\Tr A-\Tr[\sqrt A\sqrt B]$ and
$\|\Xi_0\|_2^2=\Tr A+\Tr B-2\Tr[\sqrt A\sqrt B]$, and $\Tr A=\Tr B$.}
\lby{$\Tr[\sqrt A\Pinch(\sqrt B)]=\Tr[\Pinch(\sqrt A)\sqrt B]=\Tr[\sqrt A\sqrt B]$ by
$\langle1\rangle1$; $\|\Xi_0\|_2^2=\Tr[(\sqrt B-\sqrt A)^2]$ expanded; $\Pinch$ is trace preserving}
\lstep{1}{4}{\eqref{eq:jensen}.}\lby{$\langle1\rangle2$, $\langle1\rangle3$ and $\Tr A=\Tr B$}
\lstep{1}{5}{$\Tr[Y\Xi_0]=\Tr[Y\Pinch\Xi_0]=-\Tr[Y\Theta_0]$, and $|\Tr[Y\Theta_0]|\le c\Tr[\sqrt A\Theta_0]$.}
\lby{$Y$ block diagonal and $\langle1\rangle1$; $\Theta_0\ge0$ and $-c\sqrt A\le Y\le c\sqrt A$ give
$-c\Tr[\sqrt A\Theta_0]\le\Tr[Y\Theta_0]\le c\Tr[\sqrt A\Theta_0]$ (pair the operator inequality with
the positive functional $Z\mapsto\Tr[Z\Theta_0]$)}
\lqed{1}{6}\lby{$\langle1\rangle4$, $\langle1\rangle5$}

\begin{theorem}[cross-term-free tail bound]\label{thm:nocross}
Let $Q$ commute with $E$, $0\le\widetilde Q\le1$, $A=\Pinch(\widetilde Q)=X_{11}\oplus X_{22}$,
$S_Q=Q^{1/2}$, $C_Q=(1-Q)^{1/2}$, $S_A=A^{1/2}$, $C_A=(1-A)^{1/2}$,
$S=\widetilde Q^{1/2}$, $C=(1-\widetilde Q)^{1/2}$, $\Xi_0=S-S_A$, $\Xi_1=C-C_A$, and let $U$ be the
optimal (block-diagonal) Uhlmann unitary of Note~5, Theorem~2.1 for the pair $(Q,A)$, i.e.\
$\Fid(\omega_Q,\omega_A)=|\det M_A|$, $M_A:=S_QS_A+C_QUC_A$.  Put $\gamma:=\|M_A^{-1}\|$,
$T:=M_A^{-1}(S_Q\Xi_0+C_QU\Xi_1)$ and assume
\begin{equation}\label{eq:Hyp}
 \pm\tfrac12\bigl(M_A^{-1}S_Q+S_QM_A^{-*}\bigr)\le c_0\,S_A,\qquad
 \pm\tfrac12\bigl(M_A^{-1}C_QU+U^*C_QM_A^{-*}\bigr)\le c_1\,C_A,\qquad \|T\|<1 .
\tag{H}
\end{equation}
Then
\begin{equation}\label{eq:nocross}
 0\ \le\ \Phi-\Phi_W\ \le\ \Phi_{22}
 +\frac{c_0}{2}\|\Xi_0\|_2^2+\frac{c_1}{2}\|\Xi_1\|_2^2
 +\frac{\|T\|_2^2}{2\,(1-\|T\|)},
\end{equation}
where $\Phi_{22}:=-\log\Fid(\omega_{Q_{22}},\omega_{X_{22}})$ is the complement-block defect.
There is no term of order $\sqrt{\Phi}$.
\end{theorem}
\lstep{1}{1}{$\Fid(\omega_Q,\omega_{\widetilde Q})\ge|\det(S_QS+C_QUC)|$ for \emph{every} unitary $U$.}
\lby{Note~5, Lemma~3.1 of \texttt{tail\_bound.tex} with the purification
$\Theta^U_{\widetilde Q}=\binom{S}{UC}$, which is an isometry with
$(\Theta^U)^*\Theta^U=S^2+C^*U^*UC=1$ and first block-symbol $\widetilde Q$; Uhlmann's theorem gives
$\Fid=\max_U|\det(\Theta_Q^*\Theta^U_{\widetilde Q})|$ and any $U$ is a lower bound}
\lstep{1}{2}{$S_QS+C_QUC=M_A(1+T)$ and hence
$\Phi\le-\log|\det M_A|-\log|\det(1+T)|$.}
\lby{$S=S_A+\Xi_0$, $C=C_A+\Xi_1$; $\langle1\rangle1$; $M_A$ is invertible because
$|\det M_A|=\Fid(\omega_Q,\omega_A)>0$}
\lstep{1}{3}{$-\log|\det M_A|=\Phi_W+\Phi_{22}$.}
\lby{$Q$ and $A$ are both block diagonal, so $U$ may be chosen block diagonal and $M_A$ is block
diagonal; $\Fid$ is multiplicative over the tensor factorisation
$\cF(EH)\otimes\cF(E^\perp H)$ (Lemma~3.1 of \texttt{tail\_bound.tex}), and $EAE=X_{11}=E\widetilde QE$,
so the window factor is $e^{-\Phi_W}$ and the complement factor is $e^{-\Phi_{22}}$}
\lstep{1}{4}{$\log|\det(1+T)|\ge\Re\Tr T-\frac{\|T\|_2^2}{2(1-\|T\|)}$.}
\lby{$\log\det(1+T)=\Tr\log(1+T)=\Tr T+\sum_{k\ge2}\frac{(-1)^{k+1}}k\Tr T^k$ for $\|T\|<1$, and
$|\Tr T^k|\le\|T\|_2^2\|T\|^{k-2}$ for $k\ge2$, so
$|\sum_{k\ge2}|\le\frac{\|T\|_2^2}{2}\sum_{j\ge0}\|T\|^{j}$}
\lstep{1}{5}{$\Re\Tr T=\Tr[Y_0\Xi_0]+\Tr[Y_1\Xi_1]$ with
$Y_0=\frac12(M_A^{-1}S_Q+S_QM_A^{-*})$, $Y_1=\frac12(M_A^{-1}C_QU+U^*C_QM_A^{-*})$, both Hermitian and
block diagonal.}
\lby{$\Tr[M_A^{-1}S_Q\Xi_0]+\overline{\Tr[M_A^{-1}S_Q\Xi_0]}
=\Tr[(M_A^{-1}S_Q+S_QM_A^{-*})\Xi_0]$ using $\Xi_0^*=\Xi_0$; $Q$, $A$, $U$, $M_A$ all block diagonal}
\lstep{1}{6}{$|\Re\Tr T|\le\frac{c_0}2\|\Xi_0\|_2^2+\frac{c_1}2\|\Xi_1\|_2^2$.}
\lby{Lemma~\ref{lem:jensen}\eqref{eq:Ytrace} applied to $(B,A)=(\widetilde Q,\Pinch\widetilde Q)$ with
$Y=Y_0$, $c=c_0$, and to $(B,A)=(1-\widetilde Q,\Pinch(1-\widetilde Q))$ with $Y=Y_1$, $c=c_1$
(note $\Pinch(1-\widetilde Q)=1-A$ and $\sqrt{1-A}=C_A$); the hypotheses are exactly \eqref{eq:Hyp}}
\lqed{1}{7}\lby{$\langle1\rangle2$--$\langle1\rangle6$; the left inequality of \eqref{eq:nocross} is
monotonicity of $\Fid$ under restriction (\texttt{tail\_bound.tex}, Theorem~5.1$\langle1\rangle1$)}

\begin{remark}[what \eqref{eq:nocross} is worth]
All four terms are quadratic in the discarded data.  $\Phi_{22}$ and $\|\Xi_1\|_2^2$, $\|\Xi_0\|_2^2$
are $O(\tau+\nu)=O(\zeta^2\kappa_c^{-2})$ by Corollary~\ref{cor:schur} and
\texttt{tail\_bound.tex}, Theorem~4.1; $\|T\|_2^2$ is $O(\tau)$ as well.  With the measured values of
\S\ref{sec:num} ($c_1\le1.011$, but $c_0\in[18,42]$) the right-hand side is $17.8$--$19.9\,\%$ of $\Phi_W$ at $\kappa_c=25.3$ (\S\ref{sec:rcr}, (C2)),
instead of the $82$--$154\,\%$ of \eqref{eq:final} of \texttt{tail\_bound.tex} --- a factor of about $4$--$9$ ($82/19.9=4.1$ to $154/17.8=8.7$);
the $1\,\%$ level is reached only by the computable form \eqref{eq:computable} (Table~\ref{tab:rcr}).  The price is hypothesis \eqref{eq:Hyp}, an operator inequality between
computable matrices, which \S\ref{sec:num} verifies with the measured constants $c_1\le1.011$ and $c_0\in[18,42]$; the figure $\approx3\times10^{-7}$ ($0.9\,\%$ of $\Phi$) at $\zeta=1/15$, $\kappa_c=32.2$ printed here before rested on $c_0,c_1\le1.02$, which \S\ref{sec:num} does not confirm (see Correction~C2).
\end{remark}
\section{Ball arithmetic for the box symbol $Q_N$}\label{sec:ball}\label{sec:qbox}

The box symbol of Note~5, \S6 is, in the centred basis $e_j=e^{ik_j\xi}/\sqrt\Lambda$,
$k_j=2\pi j/\Lambda$, $|j|\le M=\lfloor\kappa_{\max}\Lambda/(4\pi^2)\rfloor$, $N=2M+1$,
\begin{equation}\label{eq:Qbox}
 Q^{\rm c}_{jj}=\frac12-\frac{I(k_j)}{2\pi\Lambda},\quad
 I(k)=\int_0^\Lambda(\Lambda-t)\frac{\sin kt}{\sinh(t/2)}\dd t;\qquad
 Q^{\rm c}_{jl}=-\frac{(-1)^{j-l}\bigl(C(k_j)-C(k_l)\bigr)}{2\pi\Lambda(k_j-k_l)},
\end{equation}
$C(k)=\int_0^\Lambda\frac{\cos kt-1}{\sinh(t/2)}\dd t$, followed by the absolute-basis phase
$Q\mapsto\mathrm{diag}(e^{-ik_j\xi_0})\,Q\,\mathrm{diag}(e^{ik_j\xi_0})$, $\xi_0=\log(a/L)$.
Only $2N$ one-dimensional integrals occur.  \texttt{numerics/certified/cert\_qbox.py} encloses each
of them in an \texttt{arb} ball:
\begin{itemize}[leftmargin=1.6em]
\item on $[d,\Lambda]$, $d=10^{-6}$, by \texttt{python-flint}'s \texttt{acb.integral}, i.e.\ arb's
 Petras/Gauss--Legendre algorithm with a rigorous remainder from a Cauchy estimate on a Bernstein
 ellipse (the integrand is analytic in the strip $|\operatorname{Im}t|<2\pi$);
\item on $[0,d]$ by the explicit Taylor enclosures
 $\frac{\sin kt}{\sinh(t/2)}=2k\bigl(1+e_1\bigr)$, $|e_1|\le\frac{(kd)^2}6+\frac{d^2}{24}$, and
 $\frac{\cos kt-1}{\sinh(t/2)}=-k^2t\bigl(1+e_2\bigr)$, $|e_2|\le\frac{(kd)^2}{12}+\frac{d^2}{24}$,
 which follow from $\frac{\sin x}x=1+O(x^2/6)$, $\frac{x}{\sinh x}=1+O(x^2/6)$ with alternating,
 monotonically decreasing Taylor coefficients.
\end{itemize}
No quadrature rule is needed for the off-diagonal entries beyond the two values $C(k_j)$, $C(k_l)$.
\begin{center}\small
\begin{tabular}{l r r r}\toprule
$(\Lambda,\kappa_{\max})$ & $N$ & max ball radius of $Q_N$ & $\max_{jl}|Q^{\rm ball}_{jl}-Q^{\rm trap}_{jl}|$\\\midrule
$(60,30)$ & 91 & $5.6\times10^{-18}$ & $4.7\times10^{-11}$\\
\bottomrule\end{tabular}
\end{center}
Two consequences.  (i) The certified entrywise accuracy of $Q_N$ is $5.6\times10^{-18}$ at
$200$ bits, limited only by the Taylor cut $d$ (it improves as $d^3$: the $[0,d]$ contribution is $O(k\Lambda d)$ with relative error $O((kd)^2)$); with
$\|Z\|_1\le N^{3/2}\max_{jl}|Z_{jl}|$ this gives $\eta_Q\le5\times10^{-15}$ and, by
\texttt{tail\_bound.tex}~(8.1), an arithmetic margin $2\sqrt{2\Phi}\sqrt{5\eta_Q}=2.6\times10^{-9}$ at
$\zeta=1/15$, i.e.\ $7.5\times10^{-5}$ relative --- PT's $3\,\%$ margin for $Q_N$ is reduced by a
factor $400$.  (ii) \emph{The $4\times10^5$-node trapezoid rule of \texttt{compression\_box.Q\_box} is
accurate to only $4.7\times10^{-11}$ entrywise}, and all $N^2-N$ off-diagonal float entries lie
\emph{outside} their rigorous enclosures (the $N$ diagonal ones lie inside).  This is the
Euler--Maclaurin endpoint term $\frac{h^2}{12}[f'(\Lambda)-f'(0)]$ of the trapezoid rule and had not
been quantified; it corresponds to $\eta\approx4\times10^{-8}$ in trace norm, for which
\texttt{tail\_bound.tex}~(8.1) gives arithmetic margins of $3.3\,\%$, $21.7\,\%$ and $204\,\%$ at
$\zeta=8/15$, $1/15$ and $1/120$ respectively --- the margin is \emph{not} a single percentage, and
at small $\zeta$ it exceeds the value being computed.  All numbers in \S\ref{sec:rcr} use the
ball $Q_N$.

\begin{verdictgap}
The $Q_N$ half of the arithmetic audit requested in \texttt{tail\_bound.tex}, \S8 is complete and
rigorous.  It shows the previous margin was \emph{under}-estimated (the trapezoid error is
$4.7\times10^{-11}$, not the $\le10^{-11}$ implicitly assumed), and removes it.  The defect matrix
$\widehat D$ is a genuine two-dimensional quadrature and is \emph{not} certified here; see
\S\ref{sec:dhat}.
\end{verdictgap}

\section{The defect matrix $\widehat D$}\label{sec:dhat}

$\widehat D_{jl}=\frac1\Lambda\iint e^{-ik_j\xi}\widetilde D(\xi,\xi')e^{ik_l\xi'}\dd\xi\dd\xi'$ is a
two-dimensional integral ($N^2$ of them), evaluated in \texttt{numerics/defect\_matrix.py} by a
Gauss--Legendre product rule on panels graded logarithmically towards the corner.  A per-entry
ball evaluation is out of reach ($N^2$ nested \texttt{acb.integral} calls).  Two rigorous routes
exist and are recorded here as the remaining work:
\begin{enumerate}[label=(\alph*),leftmargin=1.8em]
\item \emph{Closed form in the modular variables.}  By \texttt{check\_note5\_theory.tex}, Prop.~6.1
 the kernel is $\widetilde F(p,p')=\zeta S_pS_{p'}\bigl[\sinh M(\sinh M+2\zeta S_pS_{p'})\bigr]^{-1}$,
 $S_p=\sinh\frac p2$, $M=\frac{p+p'}2$, a function of $(p,p')$ alone.  In the variables
 $t=p+p'$, $\sigma=p/t$ the double integral becomes
 $\int_0^\infty e^{i(\kappa_j+\kappa_l)t/4\pi}t\,g_{jl}(t)\dd t$ with
 $g_{jl}(t)=\int_0^1\widetilde F(t\sigma,t(1-\sigma))e^{i(\kappa_j-\kappa_l)(\sigma-1/2)t/2\pi}\dd\sigma$
 --- one compact inner integral and one exponentially damped outer integral, both amenable to
 \texttt{acb.integral}; cost $\sim N^2$ nested integrations, feasible but not run here.
\item \emph{Bernstein-ellipse remainder.}  $\widetilde F$ is analytic in $p,p'$ in the strip
 $|\operatorname{Im}p|<\pi$, so the $n$-point Gauss--Legendre error on a panel of half-width $h$ is
 $\le\frac{64Mh}{15(1-\rho^{-2})\rho^{2n}}$ with $\rho=\frac{\pi}{h}+\sqrt{\frac{\pi^2}{h^2}-1}$
 (Trefethen 2008, Thm.~4.5).  With $h=0.125$ and $n=12$ this is below $10^{-30}$ in the outer region;
 the graded corner panels need the same estimate in $\sigma=\log p$, where the integrand is a smooth
 function of $\tanh$'s of $\sigma$.
\end{enumerate}
Empirically the rule is converged: doubling $n_{\rm gl}$ from $12$ to $24$ and halving the panel
width changes $\widehat D$ by less than $10^{-13}$ entrywise.  We therefore carry $\widehat D$ with
the \emph{estimated} tolerance $\eta_D=10^{-12}$ (trace norm $\le N^{3/2}\eta_D\approx10^{-9}$), and
mark every enclosure below with a dagger to record that this one input is validated but not proved.
\section{Verified spectra, and the evaluation of Theorem~\ref{thm:nocross}}\label{sec:eig}\label{sec:num}

\paragraph{Working around \texttt{acb\_mat.eig}.}  PZ's blocker (arb returns \texttt{nan} once the
spectral range exceeds $\sim10^{-20}$) is avoided altogether: \emph{no} eigendecomposition of the
ball matrix is performed.  All spectral objects are built from an approximate eigenbasis $U$ of the
\emph{midpoint} matrix and are then certified a posteriori:
\begin{enumerate}[label=(\roman*),leftmargin=1.8em]
\item \emph{Eigenvalue enclosures.}  With $R:=U^*Q_NU$ a ball matrix, Gershgorin discs on $R$ enclose
 the eigenvalues of $Q_N$ rigorously; the off-diagonal mass of $R$ is $\le\|Q_NU-U\,\mathrm{diag}\|_2$
 and is $O(10^{-16})$, which separates every eigenvalue of $Q_N$ except inside the clusters at
 $O(e^{-\kappa_{\max}})$, where only the union of discs is needed --- and that is all the spectral
 \emph{window} requires.
\item \emph{Matrix square roots.}  For $M>0$ and an approximate $X\approx\sqrt M$, the Sylvester
 identity \eqref{eq:sylv} gives the a posteriori bound
 $\|\sqrt M-X\|_{1,2}\le\|M-X^2\|_{1,2}\big/\bigl(\sqrt{\lambda_{\min}(M)}+\lambda_{\min}(X)\bigr)$,
 and $|\Tr\log(1+\sqrt M)-\Tr\log(1+X)|\le\|\sqrt M-X\|_1$ because $\log(1+\cdot)$ is $1$-Lipschitz.
 The same identity that answers PT's question thus also certifies the numerics.
\end{enumerate}

\paragraph{The computable form of Theorem~\ref{thm:nocross}.}  Steps
$\langle1\rangle2$--$\langle1\rangle3$ of Theorem~\ref{thm:nocross} give, \emph{without any
hypothesis},
\begin{equation}\label{eq:computable}
 \boxed{\ 0\ \le\ \Phi-\Phi_W\ \le\ \Phi_{22}\ +\ \Phi_{\rm off},\qquad
 \Phi_{\rm off}:=-\log\bigl|\det(1+T)\bigr|,\ \ T=M_A^{-1}\bigl(S_Q\Xi_0+C_QU\Xi_1\bigr),\ }
\end{equation}
in which every object is an explicit function of $Q_N$ and $\widetilde Q_N$.  Hypothesis \eqref{eq:Hyp}
and Lemma~\ref{lem:jensen} are what \emph{guarantee} that $\Phi_{\rm off}$ is quadratic; the numbers
below evaluate it.  Script: \texttt{numerics/certified/cert\_tail\_np.py} (double precision for
$\Xi_0,\Xi_1,T$, \texttt{mpmath} at $60$ digits for $\Phi_{22}$; $\Phi_W$ from
\texttt{results\_A2.txt}/\texttt{results\_hp2.txt}).  The double-precision $\Phi_{\rm off}$ entries below are unreliable (float64 noise, \S\ref{sec:rcr}, item (P)), so the widths and gains computed from them are withdrawn; at $\kappa_c=25.3$ the corrected $60$-digit values are in Table~\ref{tab:rcr} (see Correction~C2).

\begin{center}\small
\begin{tabular}{r rr rr rr}\toprule
&\multicolumn{3}{c}{$\kappa_c=18.4$}&\multicolumn{3}{c}{$\kappa_c=25.3$}\\
\cmidrule(lr){2-4}\cmidrule(lr){5-7}
$\zeta$&$\Phi_{22}$&$\Phi_{\rm off}$&width&$\Phi_{22}$&$\Phi_{\rm off}$&width\\\midrule
$1/120$&$5.46(-9)$&$6.45(-9)$&$2.10\%$&$1.94(-9)$&$2.60(-9)$&$0.79\%$\\
$1/60$ &$2.64(-8)$&$4.61(-8)$&$3.22\%$&$8.00(-9)$&$1.91(-8)$&$1.19\%$\\
$1/30$ &$1.19(-7)$&$1.42(-7)$&$2.95\%$&$3.26(-8)$&$2.46(-8)$&$0.64\%$\\
$1/15$ &$5.14(-7)$&$5.44(-7)$&$3.08\%$&$1.31(-7)$&$1.32(-7)$&$0.76\%$\\
$2/15$ &$2.14(-6)$&$2.05(-6)$&$3.25\%$&$5.27(-7)$&$4.94(-7)$&$0.78\%$\\
$4/15$ &$8.73(-6)$&$7.77(-6)$&$3.60\%$&$2.11(-6)$&$1.68(-6)$&$0.81\%$\\
$8/15$ &$3.49(-5)$&$3.03(-5)$&$4.38\%$&$8.38(-6)$&$6.80(-6)$&$1.00\%$\\
\bottomrule\end{tabular}
\end{center}
``width'' is $(\Phi_{22}+\Phi_{\rm off})/\Phi_W$; the corresponding widths from \eqref{eq:final} of
\texttt{tail\_bound.tex} are $82$--$154\,\%$, so the gain is a factor $50$ at $\kappa_c=18.4$ and
$84$--$159$ at $\kappa_c=25.3$.  Both terms scale like $\kappa_c^{-2}$ (ratio $18.4^{-2}/25.3^{-2}=1.9$,
observed $2.8$--$5.5$, a range that includes the withdrawn double-precision $\Phi_{\rm off}$; the $60$-digit $\Phi_{22}$ alone gives $2.81$--$4.16$), confirming the exponent upgrade.

\paragraph{Hypothesis \eqref{eq:Hyp}.}  The measured constants are $c_1\le1.011$ (hole side, as
expected: $\sqrt{1-Q}\simeq\sqrt{1-A}$ throughout) but $c_0\in[18,42]$ (particle side).  The large
$c_0$ comes from the \emph{window edge}, where $Q$'s eigenvalue is $\approx\eps$ while the
corresponding eigenvalue of $A=\Pinch\widetilde Q$ is smaller by up to three orders of magnitude; it
does \emph{not} come from the ultraviolet, where $A\gg Q$.  With $c_0\approx35$ the a priori bound
$\Phi_{22}+\frac{c_0}2\|\Xi_0\|_2^2+\frac{c_1}2\|\Xi_1\|_2^2+\frac{\|T\|_2^2}{2(1-\|T\|)}$ is
$20$--$30\,\%$ of $\Phi_W$ (with the measured $c_0$ of each $\zeta$ the referee's rerun gives $17.8$--$19.9\,\%$, i.e.\ $\Phi^{\rm box}\le1.20\,\Phi_W$, \S\ref{sec:rcr}, (C2)): still a factor $4$--$6$ better than \eqref{eq:final} (about $4$--$9$ for $17.8$--$19.9\,\%$), but the sharp
statement is \eqref{eq:computable}.
\section{Certified enclosures}\label{sec:table}

\begin{table}[ht]\centering\small
\caption{Certified enclosures of the \emph{box} value $\Phi^{\rm box}(\zeta)$ (units $10^{-5}$),
$\Lambda=60$, $\kappa_{\max}=30$, $\eps=10^{-11}$ ($\kappa_c=25.3$).  Lower end: $\Phi_W$, a rigorous
lower bound for $\Phi$ (compression).  Upper end: $\Phi_W+\Phi_{22}+\Phi_{\rm off}$ by
\eqref{eq:computable}.  \textbf{Superseded by \S\ref{sec:rcr}}: the $\Phi_{\rm off}$
entries used here are double precision and are numerical noise (referee report, agent RCr).  Last two
columns: adding PT's box-cutoff term $0.0641\,\zeta^2\kappa_{\max,\rm eff}^{-2}$,
$\kappa_{\max,\rm eff}=29.6$ --- an \emph{unproved asymptotic} (the coefficient $0.513/8$ of
\texttt{referee\_quadratic\_limit.tex}, Prop.~7.1, is the $\kappa_{\max}\to\infty$ limit of
$g_B-g_B^{(\kappa_{\max})}$, with no proved remainder at finite $\kappa_{\max}$), and
the ``best'' extrapolation of \texttt{tail\_bound.tex}, Table~1.  $\dagger$: the quadrature of
$\widehat D$ is validated, not proved (\S\ref{sec:dhat}).}\label{tab:cert}
\begin{tabular}{r c r c r}\toprule
$\zeta$ & certified $\Phi^{\rm box}$ ($\dagger$) & width & $\Phi$ incl.\ box cutoff & Note~5/PT best\\\midrule
$1/120$ & $[0.057444,\ 0.057897]$ & $0.79\%$ & $[0.057444,\ 0.058405]$ & $0.05816$\\
$1/60$  & $[0.227891,\ 0.230596]$ & $1.19\%$ & $[0.227891,\ 0.232628]$ & $0.23074$\\
$1/30$  & $[0.896787,\ 0.902504]$ & $0.64\%$ & $[0.896787,\ 0.910633]$ & $0.90817$\\
$1/15$  & $[3.47372,\ 3.49999]$   & $0.76\%$ & $[3.47372,\ 3.53250]$   & $3.51910$\\
$2/15$  & $[13.05904,\ 13.16113]$ & $0.78\%$ & $[13.05904,\ 13.28118]$ & $13.2393$\\
$4/15$  & $[46.51575,\ 46.89415]$ & $0.81\%$ & $[46.51575,\ 47.42576]$ & $47.2564$\\
$8/15$  & $[151.4878,\ 153.0065]$ & $1.00\%$ & $[151.4878,\ 155.0875]$ & $154.450$\\
\bottomrule\end{tabular}\end{table}

\begin{verdictok}
The brackets of Table~\ref{tab:cert} for the box value are $0.6$--$1.2\,\%$ wide (superseded: the certified widths are $0.708$--$1.010\,\%$, Table~\ref{tab:rcr}, \S\ref{sec:rcr}; see Correction~C2), and $1.54$--$2.38\,\%$ wide (Correction~C4; superseded as well, since the same double-precision $\Phi_{\rm off}$ enters; with the upper ends of Table~\ref{tab:rcr} the same box-cutoff term gives $1.59$--$2.38\,\%$) for
$\Phi$ itself once PT's box-cutoff estimate is added (extrapolated, not certified: the box-cutoff term is an unproved asymptotic, \S\ref{sec:rcr}, (C4)); the previous widths, those of the uncertified bars of \texttt{tail\_bound.tex} (Correction~C5), were
$82$--$154\,\%$.  Every ``best'' value of \texttt{tail\_bound.tex}, Table~1 lies inside the
corresponding enclosure of the last column, which is an independent confirmation of the tail
correction $\zeta^2/(15\kappa_c^2)$.  Two inputs are not yet proved: the quadrature error of
$\widehat D$ (\S\ref{sec:dhat}, estimated $10^{-12}$ entrywise, effect $\le0.1\,\%$) and the box
cutoff itself (\texttt{tail\_bound.tex}, \S7b), whose asymptotic coefficient is unproved at finite
$\kappa_{\max}$; it is the dominant uncertainty.  See \S\ref{sec:rcr} for the corrected numbers.
\end{verdictok}
\section{Summary of findings}\label{sec:findings}
\begin{enumerate}[leftmargin=1.6em]
\item The quadratic bound asked for in \texttt{tail\_bound.tex} is \emph{false} unweighted
 (Prop.~\ref{prop:cex}, blow-up $\eps^{-1}$) and \emph{true} with the Daleckii--Krein weight, in the
 sharp form of the exact identity \eqref{eq:xiexact}; positivity (Schur complement) makes the
 weighted norm $\le\tau$ (Cor.~\ref{cor:schur}), so the weight is free.
\item The Bures cross term can be removed: Theorem~\ref{thm:nocross}/\eqref{eq:computable} replaces
 $0.173\zeta^2\kappa_c^{-1}$ by a computable quantity of order $\zeta^2\kappa_c^{-2}$, using that the
 optimal Uhlmann unitary of the pinched pair is block diagonal and that $\Pinch\sqrt{\widetilde Q}\le
 \sqrt{\Pinch\widetilde Q}$ (Lemma~\ref{lem:jensen}).
\item $Q_N$ is now enclosed to $5.6\times10^{-18}$ per entry in ball arithmetic from $2N$ certified
 one-dimensional integrals; the trapezoid rule previously used is accurate to only
 $4.7\times10^{-11}$, which gives arithmetic margins of $3.3\,\%$, $21.7\,\%$ and $204\,\%$ at $\zeta=8/15$, $1/15$, $1/120$ (\S\ref{sec:qbox}): not a single percentage, and not PT's three percent (see Correction~C2).
\item Certified brackets: \S\ref{sec:rcr}, Table~\ref{tab:rcr} (Table~\ref{tab:cert} is superseded).
\end{enumerate}
% 2026-09-24: the end-of-document command that stood on this line was moved to the end of the file, so that the RCr section below is typeset (Correction C1).

\section{Corrections after the referee report (RCr)}\label{sec:rcr}

The referee report \texttt{referee\_certified\_numerics.tex} raises one proof-breaking point and five
repairs.  All are accepted; this section records them and gives the corrected numbers.  Repairs to the
body have been made in place: \S\ref{sec:scope} no longer advertises a constant $C_\star$, verified
Rayleigh--Ritz enclosures or ball determinants (none of which is delivered);
Corollary~\ref{cor:schur} now carries the hypothesis $[X_{22},E^\pm]=0$, the correct weight
$(1-\tau_-)^{-1}$ on the $E^-$ side and the explicit use of $\widetilde Q\le1$ as well as
$\widetilde Q\ge0$; \S\ref{sec:qbox} no longer claims a single ``$8\,\%$'' margin; and the box-cutoff
term is now labelled an unproved asymptotic.

\paragraph{(P) The $\Phi_{\rm off}$ column of Table~\ref{tab:cert} was float64 noise.}  The referee
finds that under $10^{-16}$ input jitter $\Phi_{\rm off}(1/120)$ ranges over
$[-1.2\times10^{-10},\,2.9\times10^{-9}]$, and that the control identity
\begin{equation}\label{eq:ctrl}
 -\log\bigl|\det M_A\bigr|\ =\ \Phi_W+\Phi_{22}
\end{equation}
--- which is step $\langle1\rangle3$ of Theorem~\ref{thm:nocross} and must hold exactly --- failed by
$1.0\times10^{-4}$ in the double-precision run, $180$ times the quantity being computed.  This is
correct and the column is withdrawn.  The identity itself is true: the referee verifies it in
\texttt{mpmath} at $60$ digits ($1.9359858\times10^{-9}$ against $\Phi_{22}=1.935986\times10^{-9}$),
and so does the high-precision run below.

\paragraph{The high-precision path.}  \texttt{numerics/certified/cert\_tail\_hp.py} recomputes
everything in \texttt{mpmath} at $\mathtt{dps}=60$ from the ball-certified $Q_N$ of
\S\ref{sec:qbox} (midpoints; the ball radius $5.62\times10^{-18}$ per entry is carried separately),
with no double-precision object entering any spectral function:
\begin{itemize}[leftmargin=1.6em]
\item the eigenbasis of $Q_N$ is computed by \texttt{mp.eighe} at $60$ digits (a float eigenbasis is
 used only as a pre-rotation, which is an exact unitary change of basis and cannot affect the result);
\item $\sqrt{\widetilde Q}$, $\sqrt{1-\widetilde Q}$ come from one $60$-digit eigendecomposition of
 $\widetilde Q$; $\sqrt A$, $\sqrt{1-A}$, $G_A^{1/2}$ and the polar unitary $U$ are computed
 \emph{blockwise} (the pinched matrix is block diagonal), which is both faster and exact;
\item $\Phi_W$ and $\Phi_{22}$ are evaluated by Note~5, Theorem~2.1 at $60$ digits, so that
 $\Phi_W$ no longer comes from \texttt{results\_hp2.txt} (repair (2)); the shift against the
 trapezoid-based value is recorded in Table~\ref{tab:rcr};
\item $\log|\det(\cdot)|$ is evaluated by a $60$-digit LU factorisation with partial pivoting, and the
 control identity \eqref{eq:ctrl} is printed for every run as an a posteriori check.
\end{itemize}

\paragraph{Repair (2): provenance of $\Phi_W$.}  In Table~\ref{tab:cert} the lower end was taken from
\texttt{results\_hp2.txt}, which is built on the \emph{trapezoid} $Q_N$; this contradicted the claim
that all numbers use the ball $Q_N$.  In Table~\ref{tab:rcr} $\Phi_W$ is recomputed at $60$ digits
from the ball midpoints.  The relative shift is $8.6\times10^{-6}$ and is in the safe direction (the
trapezoid value is the smaller of the two, so it was still a valid lower end); it is nonetheless well
above the $7.5\times10^{-5}$ arithmetic margin claimed for $Q_N$ in \S\ref{sec:qbox} only because the
latter is a bound on the \emph{fidelity}, not on $\Phi_W$ itself, whose relative sensitivity is larger
by $1/\Phi_W$.  Both values are quoted.

\paragraph{Repair (6): what is certified, and what is not.}  With the corrections above, the
statements this note is entitled to make are, in decreasing order of strength:
\begin{enumerate}[label=(C\arabic*),leftmargin=2.4em]
\item \emph{Fully rigorous, no unproved input}: the entrywise enclosure of the box symbol,
 $|Q^{\rm ball}_{jl}-Q_{jl}|\le5.62\times10^{-18}$ at $(\Lambda,\kappa_{\max})=(60,30)$,
 $200$ bits (\S\ref{sec:qbox}); and, as a by-product, that the $4\times10^5$-node trapezoid rule of
 \texttt{compression\_box.Q\_box} is in error by up to $4.7\times10^{-11}$ per entry.
\item \emph{Rigorous modulo the quadrature of $\widehat D$} (\S\ref{sec:dhat}) \emph{and the finite
 dimensionality of the box}: $\Phi_W\le\Phi^{\rm box}$, and the a priori bound of
 Theorem~\ref{thm:nocross} with the measured constants $c_0,c_1$, which gives
 $\Phi^{\rm box}\le1.20\,\Phi_W$ at $\kappa_c=25.3$, i.e.\ widths $17.8$--$19.9\,\%$ over
 $\zeta=1/120,\dots,8/15$.
\item \emph{Same, using the computable form \eqref{eq:computable} evaluated at $60$ digits}:
 Table~\ref{tab:rcr}.  This is the statement that reaches the $1\,\%$ level; it rests on the
 $60$-digit evaluation passing the control identity \eqref{eq:ctrl}, which is printed for each run.
\item \emph{Not certified}: any statement about $\Phi$ itself rather than $\Phi^{\rm box}$, because
 the box-cutoff term is an unproved asymptotic; and the $\widehat D$ quadrature.
\end{enumerate}

\paragraph{The corrected numbers.}  \texttt{numerics/certified/hp.log}, $\mathtt{dps}=60$,
$(\Lambda,\kappa_{\max})=(60,30)$, $\eps=10^{-11}$ ($\kappa_c=25.33$, $n_{\rm win}=81$ of $N=91$),
$\approx200$\,s per $\zeta$.  The control identity \eqref{eq:ctrl} is satisfied with residual
$|{-\log|\det M_A|}-\Phi_W-\Phi_{22}|\le1.5\times10^{-26}$ in every run (double precision:
$1.0\times10^{-4}$), i.e.\ $17$ orders of magnitude below the quantity computed.

\begin{table}[ht]\centering\small
\caption{Corrected certified enclosures of the box value $\Phi^{\rm box}(\zeta)$ (units $10^{-5}$),
$\mathtt{dps}=60$, from the ball-certified $Q_N$.  $\Phi_W$ recomputed at $60$ digits (relative shift
$+8.7\times10^{-6}$ against \texttt{results\_hp2.txt}, i.e.\ the trapezoid value was smaller and
therefore still a valid lower end).  Upper end $=\Phi_W+\Phi_{22}+\Phi_{\rm off}$ by
\eqref{eq:computable}.  $\dagger$: conditional on the $\widehat D$ quadrature (\S\ref{sec:dhat}).
This table supersedes Table~\ref{tab:cert}.}\label{tab:rcr}
\begin{tabular}{r l l l c}\toprule
$\zeta$ & $\Phi_{22}$ & $\Phi_{\rm off}$ & certified $\Phi^{\rm box}$ ($\dagger$) & width\\\midrule
$1/120$ & $1.9359517(-9)$ & $2.1285135(-9)$ & $[0.05744436,\ 0.05785080]$ & $0.708\%$\\
$1/60$  & $8.0023488(-9)$ & $8.2494114(-9)$ & $[0.22789299,\ 0.22951817]$ & $0.713\%$\\
$1/30$  & $3.2547619(-8)$ & $3.2216241(-8)$ & $[0.89679497,\ 0.90327136]$ & $0.722\%$\\
$1/15$  & $1.3127838(-7)$ & $1.2567975(-7)$ & $[3.4737509,\ 3.4994467]$   & $0.740\%$\\
$2/15$  & $5.2709911(-7)$ & $4.8610897(-7)$ & $[13.059159,\ 13.160480]$   & $0.776\%$\\
$4/15$  & $2.1089359(-6)$ & $1.8524162(-6)$ & $[46.516220,\ 46.912355]$   & $0.852\%$\\
$8/15$  & $8.3833203(-6)$ & $6.9171541(-6)$ & $[151.48958,\ 153.01963]$   & $1.010\%$\\
\bottomrule\end{tabular}\end{table}

Three things change relative to Table~\ref{tab:cert}.  (i) The widths are now a smooth, monotone
function of $\zeta$, $0.708\,\%\to1.010\,\%$, instead of the scattered $0.64$--$1.19\,\%$ of the
double-precision run --- the scatter was the noise the referee identified.  (ii) $\Phi_{\rm off}$ is
\emph{not} systematically smaller than the double-precision value of the table in \S\ref{sec:num} ($\kappa_c=25.3$): it is smaller at $\zeta=1/120$, $1/60$, $1/15$, $2/15$ (e.g.\ $2.13\times10^{-9}$ against
$2.60\times10^{-9}$ at $\zeta=1/120$) but larger at $\zeta=1/30$, $4/15$, $8/15$ ($3.22\times10^{-8}$, $1.85\times10^{-6}$, $6.92\times10^{-6}$ against $2.46\times10^{-8}$, $1.68\times10^{-6}$, $6.80\times10^{-6}$; see Correction~C3), and it tracks $\Phi_{22}$ closely, $\Phi_{\rm off}/\Phi_{22}\in
[0.83,1.10]$, as the structure of \eqref{eq:computable} predicts (both are $O(\tau)$).  (iii) The
enclosures are narrower at $\zeta=1/120$, $1/60$, $1/15$, $2/15$, where the corrected upper ends lie below the previous ones, but wider at $\zeta=1/30$, $4/15$, $8/15$, where they lie above them (upper ends $0.90327136$, $46.912355$, $153.01963$ against $0.902504$, $46.89415$, $153.0065$; widths $0.722$, $0.852$, $1.010\,\%$ against $0.64$, $0.81$, $1.00\,\%$; Tables~\ref{tab:rcr} and~\ref{tab:cert}; see Correction~C3).

\begin{verdictok}
Certified, at $\kappa_c=25.33$ and conditional only on the $\widehat D$ quadrature: the box value
$\Phi^{\rm box}(\zeta)$ is enclosed to $0.71$--$1.01\,\%$ for $\zeta=1/120,\dots,8/15$, against
$82$--$154\,\%$ in \texttt{tail\_bound.tex}, Table~1.  Unconditionally certified: the entrywise
enclosure of $Q_N$ ($5.62\times10^{-18}$) and, from it, that the previously used trapezoid rule errs
by $4.7\times10^{-11}$ per entry.  Certified without the $60$-digit evaluation, from the a priori
form of Theorem~\ref{thm:nocross} with the measured $c_0,c_1$: $\Phi_W\le\Phi^{\rm box}\le
1.20\,\Phi_W$ (widths $17.8$--$19.9\,\%$).  \emph{Not} certified: the passage from $\Phi^{\rm box}$
to $\Phi$, which needs the box-cutoff coefficient $0.513/8$ at finite $\kappa_{\max}$ --- an
asymptotic statement only.
\end{verdictok}

\section{Corrections of 2026-09-24}\label{corr0924:sec}

These corrections were made on 2026-09-24 from the project's knowledge base; the card concerned
is \texttt{certified-numerics}.  Each corrected line above was replaced by exactly one line, so the
line numbers of the text above this section are unchanged.  The items are numbered C1, C2,
\dots; ``Correction~C$n$'' refers to them, not to the certification levels (C1)--(C4) of
\S\ref{sec:rcr}.

\paragraph{C1. The RCr section is now part of the document.}  \emph{Said:} the command
\texttt{\textbackslash end\{document\}} stood at line~561, before \S\ref{sec:rcr}, so the RCr
section, with Table~\ref{tab:rcr} and the certification list (C1)--(C4), was never typeset.  The
compiled document showed only the superseded Table~\ref{tab:cert}, and its references to
\S\ref{sec:rcr} and Table~\ref{tab:rcr} (lines 67, 75, 78, 414, 517, 544, 559) were undefined.
\emph{Now:} line~561 is a comment, and \texttt{\textbackslash end\{document\}} stands at the end of
the file, after this section.  \emph{Why:} knowledge-base card \texttt{certified-numerics}, whose
statement cites \S\ref{sec:rcr} and Table~\ref{tab:rcr} as the valid version (document error
recorded on 2026-09-23).  \emph{Compiling the RCr section}, which had never been compiled: it
compiles without errors and all its references resolve, so nothing inside it had to be changed.
The only undefined references left in the document are \texttt{eq:taudef} (line~245) and
\texttt{eq:final} (lines 368, 498, 509), labels of \texttt{tail\_bound.tex}; they were undefined
before as well.  Every entry of Table~\ref{tab:rcr} agrees with
\texttt{numerics/certified/hp.log}, and the control residual printed there is at most
$1.5\times10^{-26}$.

\paragraph{C2. Body statements that contradicted \S\ref{sec:rcr}.}  With \S\ref{sec:rcr} typeset,
eight places in the body contradicted it; (f)--(h) were added after the referee REF-DOC-SMALL.  They now agree with it, and each of (a)--(e) points here.
\begin{enumerate}[label=(\alph*),leftmargin=2.4em]
\item \emph{Scope, lines 63--64.}  \emph{Said:} the a priori route ``is four to six times weaker
 than the computable form''.  \emph{Now:} it is about four to nine times better than the
 $82$--$154\,\%$ of \texttt{tail\_bound.tex} but much weaker than the computable form:
 $17.8$--$19.9\,\%$ against $0.708$--$1.010\,\%$ of $\Phi_W$ at $\kappa_c=25.3$.  \emph{Why:} (C2)
 and Table~\ref{tab:rcr}; $82/19.9=4.1$ and $154/17.8=8.7$.  The factor $4$--$6$ quoted in
 \S\ref{sec:num} and in the referee report (\texttt{referee\_certified\_numerics.tex}, sections (c)
 and ``What may be quoted'') is too narrow for $17.8$--$19.9\,\%$ against $82$--$154\,\%$
 (referee REF-DOC-SMALL, 2026-09-24).
\item \emph{Remark after Theorem~\ref{thm:nocross}, lines 367--370.}  \emph{Said:} with the
 measured values of \S\ref{sec:num} the right-hand side of \eqref{eq:nocross} is
 $\approx3\times10^{-7}$ at $\zeta=1/15$, $\kappa_c=32.2$, a factor $90$ below the final bound of
 \texttt{tail\_bound.tex} and $0.9\,\%$ of $\Phi$, and \S\ref{sec:num} verifies \eqref{eq:Hyp}
 with $c_0,c_1\le1.02$.  \emph{Now:} \S\ref{sec:num} measures $c_1\le1.011$ but $c_0\in[18,42]$;
 with these constants the right-hand side is $17.8$--$19.9\,\%$ of $\Phi_W$ at $\kappa_c=25.3$, a
 factor of about $4$--$9$ ($82/19.9=4.1$ to $154/17.8=8.7$) below the $82$--$154\,\%$ of that final bound, and the $1\,\%$ level is reached
 only by \eqref{eq:computable}.  \emph{Why:} \S\ref{sec:num} (``$c_0\in[18,42]$ (particle
 side)'', caused by the window edge) and (C2), $\Phi^{\rm box}\le1.20\,\Phi_W$ at
 $\kappa_c=25.3$; the referee report, section (c), states that \eqref{eq:Hyp} cannot hold with
 $c_0=O(1)$.  The card states the a priori bound as $\Phi_W\le\Phi^{\rm box}\le1.20\,\Phi_W$.
\item \emph{Section~\ref{sec:num}, line 482.}  \emph{Said:} nothing about the reliability of the
 double-precision $\Phi_{\rm off}$ columns, whose widths and gains (a factor $50$, and $84$--$159$)
 the text then quotes.  \emph{Now:} these entries are marked unreliable, and the widths and gains
 built on them are withdrawn; at $\kappa_c=25.3$ the corrected values are in Table~\ref{tab:rcr}.
 \emph{Why:} \S\ref{sec:rcr}, item (P): the double-precision $\Phi_{\rm off}$ is float64 noise, and
 the column is withdrawn.
\item \emph{Verdict after Table~\ref{tab:cert}, lines 537--538.}  \emph{Said:} ``The certified
 brackets for the box value are $0.6$--$1.2\,\%$ wide, and $1.7$--$2.4\,\%$ wide for $\Phi$ itself
 once PT's box-cutoff estimate is added''.  \emph{Now:} these are the brackets of the superseded
 Table~\ref{tab:cert}, the certified widths are the $0.708$--$1.010\,\%$ of Table~\ref{tab:rcr},
 and the brackets for $\Phi$ itself are extrapolated, not certified.  \emph{Why:}
 Table~\ref{tab:rcr} and (C4): ``Not certified: any statement about $\Phi$ itself rather than
 $\Phi^{\rm box}$, because the box-cutoff term is an unproved asymptotic''.
\item \emph{Summary of findings, item 3, line 558.}  \emph{Said:} the trapezoid error
 $4.7\times10^{-11}$ is ``an eight-percent (not three-percent) arithmetic margin''.  \emph{Now:} it
 gives margins of $3.3\,\%$, $21.7\,\%$ and $204\,\%$ at $\zeta=8/15$, $1/15$, $1/120$, not a
 single percentage.  \emph{Why:} \S\ref{sec:qbox} (lines 411--414), which \S\ref{sec:rcr} records
 as repaired so that it ``no longer claims a single `$8\,\%$' margin''.
\item \emph{Section~\ref{sec:num}, line 509} (added after the referee REF-DOC-SMALL).
 \emph{Said:} the a priori bound with $c_0\approx35$ ``is $20$--$30\,\%$ of $\Phi_W$: still a factor
 $4$--$6$ better than'' the final bound of \texttt{tail\_bound.tex}.  \emph{Now:} the same, plus
 the measured values: with the measured $c_0$ of each $\zeta$ the referee's rerun gives
 $17.8$--$19.9\,\%$, i.e.\ $\Phi^{\rm box}\le1.20\,\Phi_W$, a factor of about $4$--$9$.
 \emph{Why:} (C2) of \S\ref{sec:rcr} and the referee report, section (d) (``$17.8$--$19.9\,\%$ of
 $\Phi_W$ across the seven $\zeta$ (the note says $20$--$30\,\%$)'').
\item \emph{Scope, lines 66--67} (a repeat of (c), added after the referee REF-DOC-SMALL).
 \emph{Said:} the scope ``replaces the factor-$80$ overshoot of \texttt{tail\_bound.tex} by a
 factor $50$--$160$''.  \emph{Now:} a factor $116$--$153$ at $\kappa_c=25.3$.  \emph{Why:} the old
 factor came from the withdrawn double-precision widths; the widths $82$--$154\,\%$ of
 \texttt{tail\_bound.tex} (its Table~1, per $\zeta$: $82$, $83$, $85$, $87$, $93$, $126$,
 $154\,\%$) divided by those of Table~\ref{tab:rcr} ($0.708$ \dots\ $1.010\,\%$) give $115.8$,
 $116.4$, $117.7$, $117.6$, $119.8$, $147.9$, $152.5$, i.e.\ $116$--$153$ within the rounding of
 the printed widths.
\item \emph{Section~\ref{sec:num}, line 501.}  \emph{Said:} both terms scale like
 $\kappa_c^{-2}$, ``observed $2.8$--$5.5$''.  \emph{Now:} the range is marked as including the
 withdrawn double-precision $\Phi_{\rm off}$; the $60$-digit $\Phi_{22}$ alone gives
 $2.81$--$4.16$.  \emph{Why:} the ratios of the $\Phi_{22}$ columns at $\kappa_c=18.4$ and $25.3$
 in the table of \S\ref{sec:num} (lines 489--495) are $2.81$, $3.30$, $3.65$, $3.92$, $4.06$,
 $4.14$, $4.16$; those of the $\Phi_{\rm off}$ columns ($2.41$ to $5.77$) rest on the entries
 withdrawn in (c).
\end{enumerate}

\paragraph{C3. Remarks (ii) and (iii) after Table~\ref{tab:rcr}, lines 658--659 and 661.}
\emph{Said:} (ii) ``$\Phi_{\rm off}$ is systematically \emph{smaller} than the double-precision
value''; (iii) ``The enclosures are narrower, not wider: the corrected upper ends lie below the
previous ones.''  \emph{Now:} (ii) the $60$-digit $\Phi_{\rm off}$ of Table~\ref{tab:rcr} is smaller
than the double-precision value of the table in \S\ref{sec:num} ($\kappa_c=25.3$) at
$\zeta=1/120$, $1/60$, $1/15$, $2/15$, and larger at $\zeta=1/30$, $4/15$, $8/15$:
$3.2216\times10^{-8}$, $1.8524\times10^{-6}$, $6.9172\times10^{-6}$ against $2.46\times10^{-8}$,
$1.68\times10^{-6}$, $6.80\times10^{-6}$ (rows at lines 648, 651, 652 against lines 491, 494, 495).
(iii) The enclosures of Table~\ref{tab:rcr} are narrower than those of Table~\ref{tab:cert} at
$\zeta=1/120$, $1/60$, $1/15$, $2/15$, and wider at $\zeta=1/30$, $4/15$, $8/15$, where the
corrected upper ends lie above the previous ones: $0.90327136$, $46.912355$, $153.01963$ against
$0.902504$, $46.89415$, $153.0065$, and widths $0.722$, $0.852$, $1.010\,\%$ against $0.64$,
$0.81$, $1.00\,\%$ (rows at lines 648, 651, 652 against lines 529, 532, 533).  Checked and right
as printed: remark (i), and the range $\Phi_{\rm off}/\Phi_{22}\in[0.83,1.10]$ in (ii) (the ratios
are $0.825$ to $1.100$).  \emph{Why:} the tables of this document, row by row; the entries of
Table~\ref{tab:rcr} agree with \texttt{numerics/certified/hp.log}, and the card takes
Table~\ref{tab:rcr} as the valid version.  The two false remarks were found when \S\ref{sec:rcr}
was first typeset (C1); they are not recorded on the card.

\paragraph{C4. Width of the brackets for $\Phi$ with the box-cutoff term, line 537.}
\emph{Said:} ``$1.7$--$2.4\,\%$ wide for $\Phi$ itself once PT's box-cutoff estimate is added''.
\emph{Now:} $1.54$--$2.38\,\%$.  \emph{Why:} the column ``$\Phi$ incl.\ box cutoff'' of
Table~\ref{tab:cert} (lines 527--533) gives (upper end $-$ lower end)/lower end $=1.673$, $2.079$,
$1.544$, $1.692$, $1.701$, $1.956$, $2.376\,\%$ at $\zeta=1/120,\dots,8/15$; the smallest is the
$\zeta=1/30$ row, as the referee report says (\texttt{referee\_certified\_numerics.tex}, section
(f), item (v): ``$1.5$--$2.4\,\%$, not $1.7$--$2.4\,\%$'').  The same convention reproduces the
printed box widths of Table~\ref{tab:cert} ($0.789\,\%$ at $\zeta=1/120$, printed $0.79\,\%$).
These brackets remain extrapolated, not certified (C2(d)).  They are also superseded, since the
same double-precision $\Phi_{\rm off}$ enters their upper ends (added after the referee
REF-DOC-SMALL): adding the same box-cutoff term, $0.0641\,\zeta^2\kappa_{\max,\rm eff}^{-2}$ with
$\kappa_{\max,\rm eff}=29.6$ as in the caption of Table~\ref{tab:cert}, to the upper ends of
Table~\ref{tab:rcr} gives $1.592$, $1.605$, $1.629$, $1.676$, $1.772$, $1.970$, $2.384\,\%$, i.e.\
$1.59$--$2.38\,\%$.
\emph{Flagged, not changed} (pre-existing): in the column ``$\Phi$ incl.\ box cutoff'' of
Table~\ref{tab:cert} the upper ends at $\zeta=2/15$ ($13.28118$) and $4/15$ ($47.42576$) do not
match the caption's formula, which gives $13.29119$ and $47.41440$; the other five rows match it
to the printed digits.  The extremes of the ranges above ($\zeta=1/120$, $1/30$, $8/15$) are
unaffected.

\paragraph{C5. The bars of \texttt{tail\_bound.tex} and the Note~5 table (lines 27--29 and 538).}
\emph{Said:} this note is ``aimed at the width of the certified brackets for $\Phi(\zeta)$ \dots\
of Note~5, Table~3, which \texttt{tail\_bound.tex} (agent PT) left at $82$--$154\,\%$'', and the
verdict after Table~\ref{tab:cert} called those $82$--$154\,\%$ ``the previous certified widths''.
\emph{Now:} error brackets, and the bars of \texttt{tail\_bound.tex} are said to be not certified;
the table is Note~5, Table~4 (\texttt{tab:scan}).  \emph{Why:} \texttt{tail\_bound.tex} now says
that its error bars are not certified, because their upper half rests on the numerical inputs of
its final bound (its label \texttt{eq:final}): every bar is \texttt{eq:bar}, whose upper end is
that bound at $\kappa_{\rm eff}$, as \texttt{numerics/pt\_bars.py} computes and reproduces; none of
them uses the bound $\Phi-\Phi_W\le\Phi_{22}+\Phi_{\rm off}$ of this note (\texttt{tail\_bound.tex},
its Corrections C3, C4 and C8).  In the auxiliary file of Note~5
(\texttt{fidelity\_recovered\_quasifree.aux}) the table of $\Phi(\zeta)$ (tail-corrected), label
\texttt{tab:scan}, is Table~4; \texttt{tail\_bound.tex} records the same renumbering (its
Correction C3).

\end{document}
